% #############################################################################
% RESUMO em Português
% !TEX root = ../main.tex
% #############################################################################
% reset acronyms
\acresetall
% use \noindent in firts paragraph
\noindent A adoção generalizada de modelos de deep learning tem impulsionado avanços tecnológicos significativos em domínios como saúde, produtividade de desenvolvimento, análise de dados e educação. No entanto, executar inferência em modelos pouco acedidos, como os hospedados em plataformas como Hugging Face, apresenta desafios de eficiência de custos. Alugar GPUs de alto desempenho, como H100 ou RTX 4090, para tarefas de baixo throughput é economicamente impraticável. Soluções existentes como DynamoLLM~\cite{dynamollm} reduzem custos através de escalonamento de frequência e instâncias de GPU, mas carecem de suporte para escalonamento vertical heterogéneo entre diferentes tipos de hardware.

Este trabalho aborda estas limitações propondo um servidor proxy containerizado capaz de selecionar dinamicamente recursos de hardware com base em análises de custo-benefício. A nossa abordagem melhora o throughput por dólar encaminhando pedidos para recursos mais baratos, como GPUs RTX 4090 ou instâncias CPU, para inferência de baixa demanda, enquanto muda perfeitamente para hardware de alto desempenho para satisfazer requisitos de throughput de produção. O sistema utiliza containers Docker para isolar cargas de trabalho de inferência e llama.cpp como motor de inferência, fornecendo uma API compatível com OpenAI que permite substituição direta em aplicações existentes.

Implementar comutação adaptativa entre inferência CPU e GPU apresenta desafios em orquestração de containers, monitorização de saúde e encaminhamento de cargas de trabalho. O nosso servidor proxy monitoriza métricas em tempo real incluindo tokens por segundo, taxas de pedidos e utilização de hardware para tomar decisões de escalonamento informadas. O sistema mantém pools de containers separados para diferentes configurações de hardware, permitindo transições rápidas entre tipos de recursos sem interrupção de serviço.

Ao incorporar análise de custo-benefício na camada de encaminhamento, o nosso sistema suporta escalonamento vertical---comutação dinâmica entre tipos de hardware como CPU, GPUs RTX 4090 ou H100 com base nas características da carga de trabalho. A avaliação demonstra que esta abordagem reduz o custo por token para cargas de trabalho variáveis mantendo throughput de nível de produção quando necessário. A arquitetura containerizada fornece isolamento, reprodutibilidade e integração perfeita com infraestrutura cloud existente e plataformas de orquestração.